\documentclass[UTF8,a4paper,11pt,openany]{ctexbook}
\usepackage{../common/statlearnbook}
\SLSetBookMetadata{第 5 册}{统计学习方法讲义 v2.6.0}{采样方法、主题模型与图排序}
\begin{document}
\setcounter{page}{869}
\setcounter{chapter}{35}
\setcounter{section}{3}
\setcounter{figure}{4}
\pagestyle{headings}
\chaptermark{潜在狄利克雷分配}
\sectionmark{变分 EM 推断}
\textbf{KL--ELBO 恒等式给出可计算下界。}
任意候选分布$q(h)$都把观测对数证据分成ELBO与到真实后验的非负KL散度；当变分族受限时，即使坐标更新稳定，这一间隙也未必消失。

图\ref{fig:V5-C06-elbo-geometry}同时显示长度分解与一次有限坐标更新轨迹。右侧阶梯只表达本次运行的ELBO不降，不把未知上界或局部驻点误写成全局最优。
% v2.6.0 figure UID: FIG-P689-01
% Iteration 05: use an explicit 8 mm x-shift to move the complete right-hand axis into its panel.
\tikzset{slfig-FIG-P689-01/.style={font=\fontsize{9.2pt}{11.0pt}\selectfont,every path/.append style={line cap=round,line join=round}}}
\pgfplotsset{slfig-FIG-P689-01-axis/.style={tick label style={font=\footnotesize},label style={font=\small},clip=false}}
\pgfkeysifdefined{/tikz/alt}{}{\tikzset{alt/.code={}}}
\begin{figure}[htbp]
\centering
\begin{tikzpicture}[slfig-FIG-P689-01,
  every node/.style={font=\fontsize{9.2pt}{11.0pt}\selectfont,align=center,outer sep=0pt},
  panel/.style={draw=SLRuleGray,line width=.55pt,rounded corners=3pt,
    minimum width=67mm,minimum height=53mm,inner sep=0pt},
  alt={对象—关系—结论：左侧以长度分解直接显示观测对数证据等于ELBO与非负KL间隙之和；右侧坐标更新使ELBO形成非降阶梯，但只能到坐标稳定或局部驻点，未知全局上限不应被解释为已经达到。}]

\node[panel] at (-3.85,0.10) {};
\node[font=\bfseries,text=SLDeepBlue] at (-3.85,2.42)
  {证据的长度分解};
\fill[SLMidBlue,opacity=0.08] (-6.55,0.50) rectangle (-2.55,1.25);
\fill[fill=SLAccentOrange,opacity=0.08] (-2.55,0.50) rectangle (-0.95,1.25);
\draw[draw=SLDeepBlue,line width=.95pt] (-6.55,0.50) rectangle (-0.95,1.25);
\draw[draw=SLMidBlue,line width=.90pt] (-2.55,0.50) -- (-2.55,1.25);
\node[text=SLDeepBlue] at (-4.55,0.875) {$\mathcal L(q)$：证据下界};
\node[text=SLAccentOrange] at (-1.75,0.875) {KL\ 间隙};
\draw[<->,draw=SLTextGray,line width=.58pt] (-6.55,1.46) -- (-0.95,1.46)
  node[midway,above=2.0mm] {$\log p(w)$};
\node[align=left,text width=59mm] at (-3.85,-0.15)
  {$\displaystyle \log p(w)=\mathcal L(q)+
    \operatorname{KL}\!\left(q(h)\,\|\,p(h\mid w)\right)$};
\node[align=left,text width=59mm,text=SLTextGray] at (-3.85,-1.35)
  {KL$\ge0$，故$\mathcal L(q)\le\log p(w)$；变分族受限时，间隙可保持为正。};

\node[panel] at (3.85,0.10) {};
\node[font=\bfseries,text=SLDeepBlue] at (3.85,2.42)
  {坐标更新下的 ELBO 非降阶梯};
\begin{axis}[slfig-FIG-P689-01-axis,
  at={(0.80,-1.55)},anchor=south west,xshift=8mm,yshift=-10.5mm,
  width=59mm,height=34mm,scale only axis,
  xmin=0,xmax=6.3,ymin=0,ymax=1.12,
  axis lines=left,axis line style={draw=SLTextGray,line width=.55pt},
  tick style={draw=SLRuleGray},xtick={0,1,2,3,4,5,6},ytick=\empty,
  xlabel={坐标更新轮次},xlabel style={font=\fontsize{9.0pt}{10.6pt}\selectfont},
  clip=false]
  \addplot+[const plot,draw=SLMidBlue,line width=.95pt,mark=*,mark size=1.5pt]
    coordinates {(0,0.12) (1,0.34) (2,0.50) (3,0.64) (4,0.73) (5,0.79) (6,0.80)};
  \addplot+[draw=SLAccentOrange,dashed,line width=.70pt,mark=none]
    coordinates {(0,1.02) (6.3,1.02)};
  \node[anchor=north west,font=\fontsize{9.0pt}{10.6pt}\selectfont,text=SLAccentOrange]
    at (axis cs:0.15,0.98) {未知全局上限};
  \node[anchor=north east,font=\fontsize{9.0pt}{10.6pt}\selectfont,text=SLDeepBlue]
    at (axis cs:6.18,0.48) {坐标稳定／局部驻点};
\end{axis}
\end{tikzpicture}
\captionsetup{width=.94\linewidth}
\caption{观测对数证据等于ELBO与变分分布到真实后验的KL散度之和；平均场坐标更新可使ELBO逐步不降，但非凸目标下有限运行通常只得到坐标稳定点或局部驻点，多启动比较也不构成全局最优证明}
\label{fig:V5-C06-elbo-geometry}
\end{figure}

\FloatBarrier
参数化变分分布把上述坐标更新改写为有限维优化；实际应用仍应检查多启动结果、目标值与预测表现，而不能仅凭一次稳定轨迹宣称全局收敛。
\end{document}
